% 07_math_algorithms.tex
% Complex math, algorithms, and pseudocode for Siggraph
% Overleaf compatible: YES (standard packages + algorithm2e/algorithmicx)

\documentclass[12pt,a4paper]{article}
\usepackage[utf8]{inputenc}
\usepackage[T1]{fontenc}
\usepackage{lmodern}
\usepackage[english]{babel}
\usepackage{hyperref}
\usepackage{cleveref}
\usepackage{amsmath}
\usepackage{amssymb}
\usepackage{amsthm}
\usepackage{mathtools}
\usepackage{physics}
\usepackage{algorithm2e}
\usepackage{algpseudocode}
\usepackage{listings}
\usepackage{xcolor}

% Theorem environments
\newtheorem{theorem}{Theorem}[section]
\newtheorem{lemma}[theorem]{Lemma}
\newtheorem{proposition}[theorem]{Proposition}
\newtheorem{corollary}[theorem]{Corollary}
\newtheorem{definition}[theorem]{Definition}
\newtheorem{example}[theorem]{Example}
\newtheorem{remark}[theorem]{Remark}

% Algorithm2e setup
\SetAlgoLined
\SetKwInOut{Input}{Input}
\SetKwInOut{Output}{Output}
\SetKw{Continue}{continue}
\SetKw{Break}{break}
\SetKw{Return}{return}
\SetKw{True}{true}
\SetKw{False}{false}
\SetKw{And}{and}
\SetKw{Or}{or}
\SetKw{Not}{not}

% Listings for code
\lstset{
    basicstyle=\ttfamily\small,
    keywordstyle=\color{blue},
    commentstyle=\color{green!60!black},
    stringstyle=\color{red},
    numbers=left,
    numberstyle=\tiny,
    stepnumber=1,
    numbersep=5pt,
    frame=single,
    breaklines=true,
    showstringspaces=false
}

\hypersetup{
    colorlinks=true,
    linkcolor=blue,
    citecolor=red,
    urlcolor=cyan
}

\title{Complex Math, Algorithms, and Pseudocode for Siggraph}
\author{Author Name}
\date{\today}

\begin{document}

\maketitle

\begin{abstract}
This example demonstrates advanced mathematical typesetting,
algorithm presentation, and pseudocode formatting for computer graphics
and Siggraph papers. Covers vector calculus, rendering equations,
optimization, and algorithm pseudocode.
\end{abstract}

\section{Vector and Matrix Notation}\label{sec:vectors}
Using \texttt{physics} package for clean notation:

Vectors: $\vb{v}$, $\vu{v}$, $\vb*{x}$.
Matrices: $\mqty{1 & 2 \\ 3 & 4}$.
Dot product: $\vb{a} \vdot \vb{b}$.
Cross product: $\vb{a} \vcross \vb{b}$.
Norm: $\norm{\vb{v}}$, $\norm{\vb{v}}_2$.
Gradient: $\grad f$, $\grad \vb{F}$.
Divergence: $\div \vb{F}$.
Curl: $\curl \vb{F}$.
Laplacian: $\laplacian f$.

\section{Rendering Equation}\label{sec:rendering}
The rendering equation (Kajiya 1986):

\begin{equation}\label{eq:rendering}
L_o(\vb{x}, \omega_o) = L_e(\vb{x}, \omega_o) + \int_{\Omega} f_r(\vb{x}, \omega_i, \omega_o) L_i(\vb{x}, \omega_i) \abs{\vb{n} \vdot \omega_i} \d\omega_i
\end{equation}

Where:
\begin{itemize}
    \item $L_o$: outgoing radiance
    \item $L_e$: emitted radiance
    \item $f_r$: BRDF
    \item $L_i$: incident radiance
    \item $\vb{n}$: surface normal
    \item $\Omega$: hemisphere of directions
\end{itemize}

\section{Monte Carlo Integration}\label{sec:mc}
Monte Carlo estimator for \cref{eq:rendering}:

\begin{equation}\label{eq:mc}
\langle L_o \rangle = \frac{1}{N} \sum_{k=1}^N \frac{f_r(\vb{x}, \omega_k, \omega_o) L_i(\vb{x}, \omega_k) \abs{\vb{n} \vdot \omega_k}}{p(\omega_k)}
\end{equation}

Importance sampling with PDF $p(\omega)$.

\section{Optimization}\label{sec:optimization}
Gradient descent update:

\begin{equation}\label{eq:gd}
\vb{\theta}_{t+1} = \vb{\theta}_t - \eta \grad_{\vb{\theta}} \mathcal{L}(\vb{\theta}_t)
\end{equation}

Adam optimizer:
\begin{align}
m_t &= \beta_1 m_{t-1} + (1 - \beta_1) g_t \\
v_t &= \beta_2 v_{t-1} + (1 - \beta_2) g_t^2 \\
\hat{m}_t &= m_t / (1 - \beta_1^t) \\
\hat{v}_t &= v_t / (1 - \beta_2^t) \\
\vb{\theta}_{t+1} &= \vb{\theta}_t - \eta \frac{\hat{m}_t}{\sqrt{\hat{v}_t} + \epsilon}
\end{align}

\section{BRDF Models}\label{sec:brdf}
Lambertian (diffuse):
\begin{equation}
f_r(\omega_i, \omega_o) = \frac{\rho}{\pi}
\end{equation}

Microfacet (Cook-Torrance):
\begin{equation}\label{eq:cook-torrance}
f_r(\omega_i, \omega_o) = \frac{D(\vb{h}) F(\omega_i, \vb{h}) G(\omega_i, \omega_o, \vb{h})}{4 \abs{\vb{n} \vdot \omega_i} \abs{\vb{n} \vdot \omega_o}}
\end{equation}

Where $\vb{h} = \frac{\omega_i + \omega_o}{\norm{\omega_i + \omega_o}}$.

GGX distribution:
\begin{equation}
D(\vb{h}) = \frac{\alpha^2}{\pi \left( (\vb{n} \vdot \vb{h})^2 (\alpha^2 - 1) + 1 \right)^2}
\end{equation}

Schlick Fresnel:
\begin{equation}
F(\omega_i, \vb{h}) = F_0 + (1 - F_0) (1 - \omega_i \vdot \vb{h})^5
\end{equation}

\section{Algorithm: Path Tracing}\label{sec:pathtrace}
\begin{algorithm}[htbp]
\caption{Recursive Path Tracing}\label{alg:pathtrace}
\Input{Ray $r$, depth $d$, max depth $D$}
\Output{Radiance $L$}
\If{$d \ge D$}{\Return 0}
Intersect $r$ with scene $\to$ hit point $\vb{x}$, normal $\vb{n}$, material $m$\;
\If{no hit}{\Return background color}
$L \gets m.L_e$ \tcp*{emitted radiance}
Sample $\omega_i \sim p(\omega_i | \vb{n})$ \tcp*{importance sample BRDF}
$r' \gets \text{Ray}(\vb{x}, \omega_i)$
$L_i \gets \text{PathTrace}(r', d+1, D)$
$L \gets L + m.f_r(\omega_i, r.\omega_o) \cdot L_i \cdot \frac{\abs{\vb{n} \vdot \omega_i}}{p(\omega_i)}$
\Return $L$
\end{algorithm}

\section{Algorithm: BVH Traversal}\label{sec:bvh}
\begin{algorithm}[htbp]
\caption{BVH Ray Traversal (Iterative)}\label{alg:bvh}
\Input{Ray $r$, BVH root node}
\Output{Closest hit}
stack $\gets$ empty\;
push root to stack\;
$t_{min} \gets \infty$\;
\While{stack not empty}{
    node $\gets$ pop stack\;
    \If{$r$ intersects node.bounds}{
        \If{node.isLeaf}{
            \ForEach{tri in node.primitives}{
                hit $\gets$ intersect(r, tri)\;
                \If{hit $\land$ hit.t $< t_{min}$}{
                    $t_{min} \gets$ hit.t\;
                    closest $\gets$ hit\;
                }
            }
        }\Else{
            push node.right to stack\;
            push node.left to stack\;
        }
    }
}
\Return closest
\end{algorithm}

\section{Algorithm: Denoising (SVGF)}\label{sec:svgf}
\begin{algorithm}[htbp]
\caption{Spatio-Temporal Variance-Guided Filtering}\label{alg:svgf}
\Input{Noisy image $I_t$, variance $V_t$, motion vectors $M_t$}
\Output{Denoised image $\hat{I}_t$}
\ForEach{pixel $x$}{
    history $\gets$ temporal\_reproject($I_{t-1}$, $M_t$, $x$)\;
    $V_{hist} \gets$ temporal\_reproject($V_{t-1}$, $M_t$, $x$)\;
    \If{history valid $\land$ $V_{hist} < V_t(x)$}{
        $I_t(x) \gets$ lerp($I_t(x)$, history, $\alpha$)\;
        $V_t(x) \gets \min(V_t(x), V_{hist})$\;
    }
}
\For{$i = 1$ to $k$}{
    $r_i \gets$ compute\_kernel\_radius($V_t$, $i$)\;
    $I_t \gets$ atrous\_filter($I_t$, $V_t$, $r_i$)\;
}
\Return $I_t$
\end{algorithm}

\section{Algorithm with Algorithmicx}\label{sec:algorithmicx}
\begin{algorithm}[htbp]
\caption{Gradient Descent with Line Search}\label{alg:linesearch}
\begin{algorithmic}[1]
\State $\vb{x} \gets \vb{x}_0$
\State $\alpha \gets 1.0$
\For{$k = 1$ to $K$}
    \State $\vb{g} \gets \nabla f(\vb{x})$
    \If{$\norm{\vb{g}} < \epsilon$}
        \State \Return $\vb{x}$
    \EndIf
    \State $\vb{d} \gets -\vb{g}$
    \While{$f(\vb{x} + \alpha \vb{d}) > f(\vb{x}) + c_1 \alpha \vb{g}^T \vb{d}$}
        \State $\alpha \gets \rho \alpha$ \Comment{backtrack}
    \EndWhile
    \State $\vb{x} \gets \vb{x} + \alpha \vb{d}$
\EndFor
\State \Return $\vb{x}$
\end{algorithmic}
\end{algorithm}

\section{Neural Network Notation}\label{sec:nn}
MLP with ReLU:
\begin{equation}
f(\vb{x}) = \vb{W}_L \sigma(\vb{W}_{L-1} \sigma(\dots \sigma(\vb{W}_1 \vb{x} + \vb{b}_1) \dots) + \vb{b}_{L-1}) + \vb{b}_L
\end{equation}

Positional encoding (NeRF):
\begin{equation}
\gamma(\vb{x}) = [\sin(2^0 \pi \vb{x}), \cos(2^0 \pi \vb{x}), \dots, \sin(2^{L-1} \pi \vb{x}), \cos(2^{L-1} \pi \vb{x})]
\end{equation}

\section{Code Listings}\label{sec:listings}
GLSL shader snippet:
\begin{lstlisting}[language=C, caption=GLSL fragment shader, label=lst:glsl]
#version 330 core
in vec3 Normal;
in vec3 WorldPos;
out vec4 FragColor;

uniform vec3 lightPos;
uniform vec3 viewPos;
uniform vec3 lightColor;
uniform vec3 objectColor;

void main() {
    // Ambient
    float ambientStrength = 0.1;
    vec3 ambient = ambientStrength * lightColor;

    // Diffuse
    vec3 norm = normalize(Normal);
    vec3 lightDir = normalize(lightPos - WorldPos);
    float diff = max(dot(norm, lightDir), 0.0);
    vec3 diffuse = diff * lightColor;

    // Specular
    float specularStrength = 0.5;
    vec3 viewDir = normalize(viewPos - WorldPos);
    vec3 reflectDir = reflect(-lightDir, norm);
    float spec = pow(max(dot(viewDir, reflectDir), 0.0), 32);
    vec3 specular = specularStrength * spec * lightColor;

    vec3 result = (ambient + diffuse + specular) * objectColor;
    FragColor = vec4(result, 1.0);
}
\end{lstlisting}

CUDA kernel:
\begin{lstlisting}[language=C, caption=CUDA ray-triangle intersection, label=lst:cuda]
__device__ bool intersect_triangle(
    const Ray& ray,
    const float3 v0, const float3 v1, const float3 v2,
    float& t, float& u, float& v) {

    float3 e1 = v1 - v0;
    float3 e2 = v2 - v0;
    float3 h = cross(ray.dir, e2);
    float a = dot(e1, h);

    if (fabsf(a) < 1e-8f) return false;

    float f = 1.0f / a;
    float3 s = ray.orig - v0;
    u = f * dot(s, h);
    if (u < 0.0f || u > 1.0f) return false;

    float3 q = cross(s, e1);
    v = f * dot(ray.dir, q);
    if (v < 0.0f || u + v > 1.0f) return false;

    t = f * dot(e2, q);
    return t > 1e-4f;
}
\end{lstlisting}

\section{Theorems and Proofs}\label{sec:theorems}
\begin{theorem}[Rendering Equation Solution]
The rendering equation \cref{eq:rendering} has a unique solution
$L \in L^\infty(\mathcal{M} \times S^2)$ for scenes with
$\norm{f_r}_\infty < 1$.
\end{theorem}

\begin{proof}
Define the transport operator $\mathcal{T}: L \mapsto L_e + \mathcal{K}L$
where $(\mathcal{K}L)(\vb{x}, \omega_o) = \int f_r L_i \abs{\vb{n} \vdot \omega_i} \d\omega_i$.
Since $\norm{\mathcal{K}} < 1$, Neumann series $\sum_{k=0}^\infty \mathcal{K}^k L_e$ converges.
\end{proof}

\begin{lemma}[Unbiased MC Estimator]
The estimator in \cref{eq:mc} is unbiased: $\E[\langle L_o \rangle] = L_o$.
\end{lemma}

\begin{proof}
By linearity of expectation and $p(\omega)$ as the sampling density.
\end{proof}

\section{Mathematical Symbols Reference}\label{sec:symbols}
Common symbols used in graphics papers:
\begin{center}
\begin{tabular}{ll}
$\vb{x}, \vb{y}, \vb{z}$ & Position vectors \\
$\omega_i, \omega_o$ & Incident/outgoing directions \\
$\vb{n}$ & Surface normal \\
$\vb{h}$ & Half vector \\
$L, E$ & Radiance, irradiance \\
$f_r, \rho$ & BRDF, albedo \\
$\Omega, \mathcal{M}$ & Hemisphere, scene manifold \\
$\mathcal{K}$ & Transport operator \\
$\sigma_t, \sigma_s$ & Extinction, scattering coefficients \\
$\alpha, \beta$ & Roughness, metallic parameters \\
\end{tabular}
\end{center}

\end{document}